%% =====================================================================
%%  core/config/membrete.tex  --  Membretes y pie institucional
%% ---------------------------------------------------------------------
%%  Equivalente a covers/ del AWF: cada variante de papel membretado
%%  vive en membretes/<nombre>/membrete.tex y SOLO rellena el
%%  encabezado fancyhdr consumiendo \meta{...}. Se selecciona con:
%%
%%     \usarmembrete{congreso}    membrete oficial (escudo + logo CR)
%%     \usarmembrete{despacho}    variante con logo Camara de Diputados
%%     \sinmembrete               pagina limpia (segunda hoja, anexos)
%%
%%  El PIE INSTITUCIONAL (comun a todos los membretes) va al fondo de
%%  cada pagina: filete + correo institucional, N.° de despacho y
%%  telefono (si estan definidos via \setmeta) + folio opcional:
%%
%%     \confolio / \sinfolio      "Pag. x de y" a la derecha del pie
%%     \sinpieinstitucional       pagina sin pie (p. ej. anexos)
%% =====================================================================

\RequirePackage{fancyhdr}
\RequirePackage{lastpage}

% --- Ancho de la caja del nombre (el lema se centra respecto a ella) --
\newlength{\despachoAnchoCaja}
\setlength{\despachoAnchoCaja}{7.6cm}
\newcommand{\setanchocajamembrete}[1]{\setlength{\despachoAnchoCaja}{#1}}

% --- Folio opcional (bandera; el pie lo dibuja si esta activa) --------
\newif\ifdespachofolio
\despachofoliofalse
\newcommand{\confolio}{\despachofoliotrue}
\newcommand{\sinfolio}{\despachofoliofalse}

% --- Filete bajo el membrete (bandera; 2026-10-01) ----------------------
%%  Los membretes congreso y despacho dibujan un filete bajo el encabezado.
%%  \sinfileteencabezado lo quita (el proyecto de ley: quedaba pegado a la
%%  primera línea del texto); la separación la da entonces el headsep.
\newif\ifdespachofileteencabezado
\despachofileteencabezadotrue
\newcommand{\sinfileteencabezado}{\despachofileteencabezadofalse}

% --- Pie institucional ------------------------------------------------
%%  Dos lineas: direccion arriba; correo/celular/telefono + folio abajo.
%%  Solo muestra los datos definidos: si un \setmeta falta o esta
%%  vacio, su segmento (o su linea) desaparece del pie.
\newcommand{\pieinstitucional}{%
  \begin{minipage}[t]{\textwidth}%
    {\color{colorfilete}\rule{\textwidth}{0.4pt}}\\[2.5pt]
    {\fuentedatos\fontsize{7.5}{9.5}\selectfont\color{neutral700}%
      \ifmetadef{direccion}{\meta{direccion}\\}{}%
      \ifmetadef{correo}{\meta{correo}}{}%
      \ifmetadef{celular}{\ \,\textperiodcentered\ \,Cel.~\meta{celular}}{}%
      \ifmetadef{telefono}{\ \,\textperiodcentered\ \,T.~\meta{telefono}}{}%
      \hfill
      \ifdespachofolio Pág.~\thepage~de~\pageref*{LastPage}\fi}%
  \end{minipage}%
}

% --- Estilo de pagina del despacho ------------------------------------
\fancypagestyle{despacho}{%
  \fancyhf{}%
  \renewcommand{\headrulewidth}{0pt}%
  \renewcommand{\footrulewidth}{0pt}%
  \fancyfoot[C]{\pieinstitucional}%
}
\pagestyle{despacho}

% --- Pie SOLO folio ----------------------------------------------------
%%  Sin filete ni datos del despacho; queda únicamente el «Pág. x de y»
%%  (si \confolio está activo). Para documentos que NO son del despacho
%%  individual, p. ej. los oficios de la BANCADA REGIONAL, firmados por
%%  varios parlamentarios (pedido de Edison, 2026-08-23).
\newcommand{\solofoliopie}{%
  \fancyfoot[C]{%
    \begin{minipage}[t]{\textwidth}%
      {\fuentedatos\fontsize{7.5}{9.5}\selectfont\color{neutral700}%
        \hfill\ifdespachofolio Pág.~\thepage~de~\pageref*{LastPage}\fi}%
    \end{minipage}%
  }%
}

% --- Pie INSTITUCIONAL DEL CONGRESO (proyectos de ley; 2026-10-01) --------
%%  Como el Word de las proposiciones de la Camara de Diputados: a la
%%  izquierda la web del Congreso; a la derecha su direccion y central
%%  telefonica; al borde inferior de la hoja, una franja guinda a todo el
%%  ancho. Sin folio y sin los datos personales del despacho (correo,
%%  celular): el proyecto es un documento de la institucion. Los textos son
%%  datos PUBLICOS del Congreso, ajustables con \setmeta{webcongreso} /
%%  {direccioncongreso} / {centralcongreso}.
\definecolor{grismembrete}{HTML}{7F7F7F}
\definecolor{franjacongreso}{HTML}{990000}
\RequirePackage{eso-pic}
\newcommand{\usarpiecongreso}{%
  \fancyfoot[C]{%
    \begin{minipage}[b]{\textwidth}%
      {\fuentedatos\fontsize{7.5}{9.2}\selectfont\color{tinta}%
        \ifmetadef{webcongreso}{\meta{webcongreso}}{www.congreso.gob.pe}%
        \hfill
        \begin{tabular}[b]{@{}r@{}}%
          \ifmetadef{direccioncongreso}{\meta{direccioncongreso}}{Plaza de la Constituci\'on, Av.~Abancay s/n -- Lima, Per\'u}\\
          Central Telef\'onica: \ifmetadef{centralcongreso}{\meta{centralcongreso}}{311-7777}%
        \end{tabular}}%
    \end{minipage}%
  }%
  \AddToShipoutPictureBG{%
    \AtPageLowerLeft{\color{franjacongreso}\rule{\paperwidth}{0.30cm}}}%
}

% --- Selectores --------------------------------------------------------
\newcommand{\usarmembrete}[1]{\input{membretes/#1/membrete.tex}}
\newcommand{\sinmembrete}{\fancyhead{}}
\newcommand{\sinpieinstitucional}{\fancyfoot[C]{}}
